\section{Requisiti Funzionali}

\subsection{Requisiti funzionali per Operatore e Amministratore}
\begin{itemize}
    \item \textbf{RF0 - Localizzazione:} Il sistema supporterà tre lingue: inglese (EN), italiano (IT) e tedesco (DE). Tutte le stringhe visibili all'utente saranno gestite attraverso cataloghi di traduzione esterni al codice. La lingua dell'interfaccia sarà selezionabile dall'utente tramite apposito controllo, con la preferenza persistita tra le sessioni. Il backend rileverà la lingua richiesta a ogni chiamata restituendo contenuti localizzati, inclusi quelli generati lato server come email e report; date e numeri seguiranno il formato della locale selezionata. La lingua di riferimento è l'inglese, con italiano e tedesco che ne rispecchiano fedelmente la struttura.
    
    \item \textbf{RF1 - Autenticazione e Gestione Sessione (Login e Logout):} Il sistema deve consentire sia agli operatori tecnici che agli amministratori comunali di accedere tramite credenziali locali (email e password) o tramite autenticazione Google. Una volta autenticato, l'utente viene reindirizzato alla dashboard generale. Il sistema deve inoltre consentire all'utente di terminare la sessione in modo sicuro (Logout), eliminando il token locale di accesso dal dispositivo e reindirizzando alla schermata di login.

    \item \textbf{RF2 - Visualizzazione dashboard generale:} Il sistema deve consentire a entrambe le tipologie di utente di visualizzare una dashboard generale che mostri una panoramica dei consumi energetici degli edifici pubblici. La dashboard deve presentare grafici aggregati dei consumi totali e indicatori di performance energetica. 

    \item \textbf{RF3 - Ricerca edifici e visualizzazione cartografica:} Il sistema deve consentire agli utenti di individuare gli edifici pubblici sia attraverso una maschera di ricerca avanzata a criteri multipli, sia mediante una mappa geografica interattiva. L'utente deve poter filtrare gli stabili combinando parametri quali nome, indirizzo e tipologia di edificio. La ricerca testuale restituisce l'elenco delle strutture conformi ai filtri, evidenziando nome, indirizzo, superficie, tipologia di riscaldamento e stato operativo. Parallelamente, la vista su mappa posiziona marker georeferenziati con codifica cromatica in base allo stato operativo dell'edificio (attivo, inattivo, dismesso); la selezione di un edificio (da elenco o da mappa) conduce alla relativa schermata di dettaglio.
    
    \item \textbf{RF4 - Visualizzazione dettaglio edificio:} Il sistema deve permettere agli utenti di visualizzare i dettagli completi di ogni singolo edificio selezionato. La schermata di dettaglio deve mostrare le informazioni dell'edificio (il nome, l'indirizzo, la superficie e l'anno di costruzione), i dati in tempo reale (la temperatura interna, la temperatura esterna e il consumo energetico istantaneo), i grafici storici dei consumi e delle temperature e lo stato di funzionamento dei sensori installati. L'utente deve poter selezionare diversi intervalli temporali per l'analisi dei dati storici e poterli filtrare per categoria di edificio (scuole, uffici comunali, biblioteche, impianti sportivi, ecc.).

    \item \textbf{RF5 - Gestione edifici e configurazione sensori:} Il sistema deve consentire all'operatore tecnico e all’amministratore di gestire gli edifici monitorati e configurare i parametri operativi dei sensori installati; inoltre possono aggiungere nuovi edifici inserendo tutti i dati anagrafici (nome, indirizzo, superficie, tipologia di impianto), associare i sensori specificando posizione fisica e tipologia (termometri interni, contatori, temperatura esterna), e rimuovere edifici dismessi. I sensori dovranno trasmettere le misurazioni in media ogni 90 secondi. Il sistema deve validare l'univocità del codice identificativo dell'edificio, verificare la correttezza dei dati inseriti e registrare ogni modifica con timestamp e identificativo dell'operatore che l'ha effettuata.

    \item \textbf{RF6 - Gestione notifiche di allarme:} Il sistema deve permettere all'operatore tecnico e all’amministratore di visualizzare, gestire e chiudere le notifiche di allarme generate; inoltre avranno l’accesso a un pannello dedicato in cui saranno elencate tutte le notifiche attive, ordinate in base al timestamp. Da questo pannello sarà possibile visualizzare i dettagli di ogni singolo allarme, inclusi l’edificio coinvolto e i valori rilevati. Dopo aver risolto il problema segnalato, si potrà chiudere la relativa notifica.

    \item \textbf{RF7 - Visualizzazione stato sensori:} Il sistema deve permettere all'operatore tecnico e all’amministratore di monitorare lo stato di funzionamento di tutti i sensori installati negli edifici. Si deve anche poter visualizzare un elenco completo dei termometri e dei contatori, con indicazione dello stato operativo (attivo o offline), data dell'ultimo aggiornamento dati e posizione fisica.

    \item \textbf{RF8 - Monitoraggio tempo reale e storico:} Il sistema deve consentire all'operatore tecnico e all'amministratore di monitorare in tempo reale i dati provenienti da un edificio selezionato. La visualizzazione deve aggiornarsi automaticamente ogni 90 secondi mostrando temperatura interna, temperatura esterna, consumo energetico tramite un grafico delle ultime 24 ore. Deve essere possibile anche cambiare la data di analisi da quella corrente a una passata. L'operatore deve poter attivare notifiche email per variazioni significative dei parametri monitorati.

    \item \textbf{RF9 - Configurazione soglie di allarme:} Il sistema deve permettere all'amministratore e all’operatore tecnico di configurare soglie personalizzate per gli allarmi automatici relativi a consumi anomali o malfunzionamenti; si deve poter definire valori limite (massimi e minimi) per tutti i tipi di sensori supportati. Inoltre è stato implementato un sistema che gestisce le mail di allerta per valori anomali.

\end{itemize}

\subsection{Requisiti funzionali esclusivi per l'Amministratore}
\begin{itemize}
    \item \textbf{RF10 - Esportazione dati:} Il sistema deve consentire all'amministratore comunale di esportare i dati relativi ai consumi energetici in diversi formati (CSV, Excel, PDF). L'utente deve poter selezionare l'edificio o il gruppo di edifici di interesse, specificare il periodo temporale e scegliere il formato di esportazione desiderato. Il sistema deve generare un file contenente tutti i dati richiesti con informazioni complete su temperature, consumi, timestamp e eventuali anomalie registrate nel periodo selezionato.

    \item \textbf{RF11 - Gestione utenti:} Il sistema deve consentire all'amministratore comunale di gestire gli utenti inseriti nella piattaforma. L'amministratore accede a una sezione dedicata dove può inviare un invito inserendo l'indirizzo email e il ruolo assegnato al nuovo utente.
    Il sistema genera e invia automaticamente una mail all'utente contenente un link di invito univoco. Cliccando sul link, l'utente dovrà completare la propria registrazione inserendo i dati personali (nome e cognome) e impostando la password di accesso.
    L'amministratore deve inoltre poter modificare il ruolo degli utenti esistenti, nonché disabilitare temporaneamente o eliminare l'account quando necessario.
    
    \item \textbf{RF12 - Visualizzazione analisi comparative:} Il sistema deve permettere all'amministratore di visualizzare analisi comparative tra diversi edifici o gruppi di edifici. L'amministratore deve poter selezionare fino a 3 edifici e confrontare i loro consumi energetici, le temperature medie e l'efficienza energetica . Il sistema deve presentare tabelle con indicatori di performance, evidenziando gli edifici più e meno efficienti.
\end{itemize}